\subsection{Vision Pipeline Organisation}
The vision system estimates the \textbf{position and orientation of the debris}, 
then tracks its relative motion. It combines Auricam images, 
two neural networks and a geometric calculation based on camera calibration and a metrically scaled 3D model of the debris.

\begin{center}
\begin{tikzpicture}[node distance=8mm and 7mm]
\node[block] (cam) {\textbf{Auricam D35}\\Monochrome image};
\node[block,right=of cam] (ssd) {\textbf{SSD-MobileNetV2}\\Debris bounding box};
\node[block,right=of ssd] (hr) {\textbf{HRNet-W32}\\Image keypoints};
\node[block,below=of hr] (pnp) {\textbf{SQPnP}\\Position and orientation};
\node[block,left=of pnp] (ukf) {\textbf{Adaptive UKF}\\Tracking and uncertainty};
\node[block,left=of ukf] (out) {\textbf{Time-stamped result}\\Pose and motion};
\draw[arrow](cam)--(ssd);
\draw[arrow](ssd)--(hr);
\draw[arrow](hr)--(pnp);
\draw[arrow](pnp)--(ukf);
\draw[arrow](ukf)--(out);
\end{tikzpicture}
\end{center}
{\small Initialisation sequence. During tracking, the UKF uses the observed keypoints and predicts the region of interest in the next image.}

The target rate is two optical observations per second, processed as they are acquired. The detector runs at initialisation, during full-image checks and after tracking is lost. 
Between these calls, the predicted region allows processing to focus on the debris.

\subsection{Camera: Monochrome Auricam D35, 25 mm Focal Length}
We select the optical head \textbf{Sodern Auricam D35}, including its baffle,
 to capture the outline and geometric features of the debris. Monochrome imaging avoids colour demosaicing;
 the global shutter exposes all pixels simultaneously~\cite{sodern_auricam}.

\begin{tabularx}{\linewidth}{|P{36mm}|Y|}
\hline
\textbf{Characteristic} & \textbf{Configuration and published specifications}\\\hline
Sensor & $2048\times2048$ pixels; pixel pitch of $5{,}5\,\mu\mathrm{m}$; global shutter.\\\hline
Optics & 25 mm focal length; standard aperture $f/3{,}5$; field of view of $35^\circ$ diagonally, approximately $25^\circ$ along each side.\\\hline
Frame rate & Full frame, 12-bit: 2,5 frames/s. Binning $2\times2$: 10 frames/s at $1024\times1024$.\\\hline
Interface & SpaceWire, up to 160 Mbit/s.\\\hline
Power supply & 4,6--6 V; published power consumption of 2 W.\\\hline
Integration & $80\times80\times140$ mm including baffle; mass of 420 g.\\\hline
\end{tabularx}

The 25 mm focal length combines a useful image scale with a framing margin. 
A projected dimension of 30 cm at 10 m spans approximately 136 native pixels. 
The optics and mounting are calibrated, and the focus is then adjusted for the observation range of 2 to 10 m. 
Pointing complies with the published Sun exclusion angle of $38^\circ$.

\subsection{Onboard Computing: Q8S, XE2 and Interfaces}
The \textbf{Xiphos Q8S} receives the images, prepares the network inputs and runs SQPnP, 
the UKF and data logging. The \textbf{CogniSAT-XE2} from Ubotica runs the neural networks on its Myriad X processor~\cite{xiphos_q8s,ubotica_xe2}.

\begin{tabularx}{\linewidth}{|P{29mm}|Y|Y|}
\hline
\textbf{Equipment} & \textbf{Key specifications} & \textbf{Role in Argus}\\\hline
Xiphos Q8S & Zynq UltraScale+ XCZU7EG ; 4 Cortex-A53, 2 Cortex-R5 and FPGA ; 4 GB of error-correcting RAM; 5 W typical. & Reception, cropping, data preparation, pose estimation and temporal tracking.\\\hline
CogniSAT-XE2 & Myriad X ; PC/104 form factor; 80 g typical; 2,5 W typical with a representative network and Ethernet. & Execution of SSD-MobileNetV2 and HRNet-W32; return of results to the Q8S.\\\hline
Carrier board & Power supplies, connectors and interface circuits. & Camera connection and integration of the SpaceWire interface with the Q8S FPGA.\\\hline
\end{tabularx}

Images arrive via SpaceWire; the Q8S communicates with the XE2 over Ethernet. Integration verifies the camera data rate, the ARM64 Linux host software and the loading of both networks. Tests measure the latency between exposure and output, as well as the power consumption of the complete system.

The vision system coste around 22 W, including margins and conversion losses, with an allocation of 1 kg and 2U covering the camera, boards and interfaces. These design allocations are separate from the typical power consumption of the equipment.

\subsection{Neural Networks: Detection Followed by Keypoint Localisation}
\textbf{SSD-MobileNetV2} receives an image prepared at $320\times320$ pixels. It searches for the debris and returns a bounding box with a score. This architecture is documented for onboard spacecraft pose estimation~\cite{black2021pose}. The bounding box selects the region to be analysed.

\textbf{HRNet-W32} receives this region, extracted from the original image and resized to $256\times256$ pixels. It locates the keypoints associated with the 3D model: corners, junctions or asymmetric features~\cite{fanizza2023pose}. The Q8S maps their coordinates back to the native image; SQPnP initialises the pose, and the UKF then performs tracking.

The networks are trained on the ground using synthetic and real images of the debris under varied lighting conditions. The detector also receives images without debris. Onboard, the weights remain fixed and the models run in FP16. Validation on the XE2 checks keypoint accuracy after compilation, false detections and the processing rate of the complete pipeline.

